\chapter{Tecnologías utilizadas y criterio de elección}
\label{cap:tecnologias}

Este capítulo aterriza la solución en tecnologías concretas. Su objetivo no es solo enumerar herramientas, sino explicar \textbf{para qué sirve cada una, por qué se ha elegido y qué papel desempeña dentro del sistema completo}. También resulta importante distinguir entre herramientas preexistentes reutilizadas en el proyecto y componentes desarrollados específicamente en este TFG.

% ============================================================
\section{Python como lenguaje de integración principal}
\label{sec:python}
% ============================================================

El lenguaje vertebrador del proyecto es Python. La razón principal no es solo su popularidad, sino su capacidad para actuar como lenguaje de integración entre bibliotecas de IA, servicios web, automatización y procesamiento documental. En este TFG, Python permite conectar en un mismo ecosistema el backend, el indexador, el pipeline de OpenWebUI, la integración con SharePoint, el scraping, el OCR y el servidor MCP.

Además, buena parte de las herramientas reutilizadas en el proyecto exponen su mejor soporte precisamente en este lenguaje. Entre los paquetes y bibliotecas más importantes utilizados se encuentran:

\begin{itemize}
    \item \textbf{FastAPI} y \textbf{Uvicorn} para los servicios web y la API REST.
    \item \textbf{Pydantic} para validación de datos y configuración tipada.
    \item \textbf{sentence-transformers} para generación de embeddings.
    \item \textbf{qdrant-client} para interacción con la base de datos vectorial.
    \item \textbf{msal} y clientes asociados para autenticación y acceso a Microsoft Graph.
    \item \textbf{playwright} para automatización de navegador.
    \item \textbf{trafilatura}, \textbf{readability} y \textbf{BeautifulSoup} para extracción de contenido web.
    \item \textbf{PaddleOCR} para OCR de documentos escaneados.
    \item \textbf{duckduckgo\_search} para búsqueda web controlada.
    \item \textbf{fastmcp} para el servidor MCP del BOE.
    \item \textbf{prometheus-client} para exponer métricas del backend.
\end{itemize}

Esta elección también ayuda a explicar la naturaleza del trabajo realizado: el proyecto no consiste en inventar estas bibliotecas, sino en \textbf{utilizarlas de forma coordinada y construir sobre ellas la lógica que faltaba para el caso de uso corporativo}.

% ============================================================
\section{OpenWebUI: interfaz conversacional reutilizada}
\label{sec:openwebui}
% ============================================================

OpenWebUI\footnote{\url{https://openwebui.com/}} es la interfaz conversacional utilizada como punto de entrada del sistema. Es una herramienta externa, ya existente, que he reutilizado para no invertir el esfuerzo del TFG en construir una interfaz de chat desde cero. Esta decisión permite concentrar el trabajo en la parte realmente diferencial del proyecto: la integración con las fuentes documentales y la lógica del asistente.

\subsection{Arquitectura interna de OpenWebUI}

OpenWebUI se estructura como una aplicación web con tres bloques principales:

\begin{itemize}
    \item \textbf{Frontend}: interfaz tipo chat accesible desde navegador, con historial de conversaciones, carga de adjuntos y streaming de respuestas.
    \item \textbf{Backend}: lógica de autenticación, persistencia de chats y gestión de modelos.
    \item \textbf{Base de datos}: almacenamiento de usuarios, conversaciones y metadatos.
\end{itemize}

Para el lector, lo importante es entender su papel dentro del proyecto: OpenWebUI aporta la experiencia de usuario, pero no constituye la inteligencia principal del sistema. Esa inteligencia se añade mediante el pipeline desarrollado en este TFG.

\subsection{Por qué OpenWebUI}

La elección de OpenWebUI se justifica porque ya ofrece varias capacidades deseables en un entorno corporativo: interfaz madura, persistencia de conversaciones, compatibilidad con modelos locales y soporte de autenticación. Frente a hacer una interfaz propia, reduce tiempo de desarrollo y riesgo técnico.

\begin{table}[H]
\centering
\caption{Comparativa de interfaces de usuario para LLMs locales.}
\label{tab:openwebui-comparativa}
\begin{tabular}{lcccc}
\toprule
\textbf{Criterio} & \textbf{OpenWebUI} & \textbf{Gradio} & \textbf{Streamlit} & \textbf{ChatUI} \\
\midrule
UI tipo ChatGPT              & \checkmark & Parcial   & Parcial  & \checkmark \\
Compatibilidad con Ollama    & Nativa     & Manual    & Manual   & Nativa \\
Sistema de Pipelines/Plugins & \checkmark & $\times$  & $\times$ & $\times$ \\
Gestión de usuarios          & \checkmark & $\times$  & $\times$ & \checkmark \\
Persistencia de chats        & \checkmark & $\times$  & $\times$ & \checkmark \\
SSO / OIDC                   & \checkmark & $\times$  & $\times$ & $\times$ \\
Código abierto               & \checkmark & \checkmark & \checkmark & \checkmark \\
Carga de archivos adjuntos   & \checkmark & \checkmark & \checkmark & Parcial \\
\bottomrule
\end{tabular}
\end{table}

El factor decisivo es el sistema de Pipelines, porque permite insertar lógica personalizada sin modificar el núcleo de la herramienta. Gracias a ello, la interfaz sigue siendo una pieza reutilizada, pero el comportamiento final del asistente pasa a estar controlado por desarrollo propio.

\subsection{Sistema de Pipelines de OpenWebUI}
\label{subsec:pipelines}

El sistema de Pipelines es el mecanismo que convierte una interfaz genérica en un asistente especializado. Funciona como una capa intermedia entre el usuario y el modelo: recibe el mensaje, decide qué hacer con él y devuelve la respuesta final.

Desde el punto de vista del proyecto, esta capacidad es clave porque permite:

\begin{itemize}
    \item interceptar cada consulta antes de enviarla al modelo,
    \item aplicar lógica de enrutamiento,
    \item conectar con el backend propio,
    \item y conservar la interfaz original sin reescribirla.
\end{itemize}

Este es uno de los mejores ejemplos de la filosofía del TFG: reutilizar una herramienta existente cuando ya resuelve bien una capa del problema y centrar el desarrollo propio en la parte que aporta valor diferencial.

% ============================================================
\section{Ollama: ejecución local de modelos}
\label{sec:ollama}
% ============================================================

Ollama\footnote{\url{https://ollama.com/}} es la herramienta externa utilizada para ejecutar modelos de lenguaje localmente. Su función es cargar modelos, servirlos por API y gestionar su uso en el hardware disponible \cite{ollama2024docs}.

\subsection{Características técnicas}

Se ha elegido Ollama por varias razones prácticas:

\begin{itemize}
    \item \textbf{Facilita el despliegue de modelos locales}: evita construir una capa de inferencia propia desde cero.
    \item \textbf{Gestiona la VRAM de manera operativa}: carga y descarga modelos según sea necesario.
    \item \textbf{Expone una API sencilla}: esto simplifica mucho su integración con el resto del sistema.
    \item \textbf{Encaja bien en contenedores}: puede desplegarse como un servicio más dentro de Docker Compose.
\end{itemize}

\subsection{Modelos desplegados}

En el sistema se emplean varias familias de modelos ya existentes, cada una con una función concreta:

\begin{table}[H]
\centering
\caption{Modelos de IA desplegados en Ollama y su función en el sistema.}
\label{tab:modelos-ollama}
\begin{tabular}{llcc}
\toprule
\textbf{Modelo} & \textbf{Función} & \textbf{Cuantización} & \textbf{VRAM} \\
\midrule
LLaMA 3.1 8B Instruct & Chat general y RAG & Q8\_0 & $\sim$8 GB \\
LLaMA 3.1 8B Instruct & Resúmenes rápidos & Q4\_0 & $\sim$4 GB \\
Qwen 2.5 (LoRA FT)    & Variante adaptada al dominio & --- & $\sim$5 GB \\
Qwen 2.5 VL 7B        & Análisis de imágenes & --- & $\sim$5 GB \\
LLaVA 13B             & Vía histórica para visión/OCR & --- & $\sim$8 GB \\
\bottomrule
\end{tabular}
\end{table}

La tabla refuerza una idea importante: los modelos no se han creado en este TFG, pero sí se han \textbf{seleccionado, desplegado, probado y asignado a tareas concretas} dentro de una arquitectura coherente.

% ============================================================
\section{Qdrant: base de datos vectorial}
\label{sec:qdrant}
% ============================================================

Qdrant\footnote{\url{https://qdrant.tech/}} es la base de datos vectorial empleada como almacén principal del conocimiento documental del sistema \cite{qdrant2024docs}. Es una herramienta externa, pero su estructura, colecciones y criterios de uso se definen dentro del proyecto.

\subsection{Características utilizadas}

\begin{itemize}
    \item \textbf{Índice HNSW}: permite búsquedas vectoriales rápidas sobre grandes volúmenes de fragmentos.
    \item \textbf{Búsqueda híbrida}: combina similitud semántica y coincidencia léxica exacta.
    \item \textbf{Filtrado por metadatos}: habilita segmentación por fuente, documento o departamento.
    \item \textbf{Multi-colección}: permite separar documentos por contexto de uso.
\end{itemize}

Para el proyecto, esto significa que Qdrant no es solo un repositorio técnico, sino el componente que hace posible que el sistema encuentre evidencia relevante antes de responder.

% ============================================================
\section{FastAPI: backend del sistema}
\label{sec:fastapi}
% ============================================================

FastAPI\footnote{\url{https://fastapi.tiangolo.com/}} es el framework seleccionado para construir los servicios backend en Python \cite{fastapi2024docs}. A diferencia de OpenWebUI, Ollama o Qdrant, que son herramientas reutilizadas, aquí sí aparece una parte importante de desarrollo propio: los servicios implementados con FastAPI contienen gran parte de la lógica central del sistema.

Se ha elegido por tres motivos principales:

\begin{itemize}
    \item \textbf{Facilidad para construir APIs tipadas y mantenibles}.
    \item \textbf{Soporte asíncrono}: útil cuando el sistema depende de múltiples llamadas a servicios externos o internos.
    \item \textbf{Integración natural con el ecosistema Python}: permite usar en el mismo entorno paquetes de IA, scraping, OCR y autenticación.
\end{itemize}

En este TFG, FastAPI sustenta el backend RAG, el indexador de documentos y parte de las integraciones auxiliares.

% ============================================================
\section{LiteLLM: proxy unificado de modelos}
\label{sec:litellm}
% ============================================================

LiteLLM\footnote{\url{https://docs.litellm.ai/}} es una herramienta externa que se utiliza como capa intermedia entre los servicios propios y los modelos servidos por Ollama \cite{litellm2024docs}. Su utilidad está en unificar llamadas, aplicar alias de modelos, gestionar caché y simplificar el enrutamiento técnico.

Desde la perspectiva del lector no técnico, su papel puede resumirse así: en lugar de que cada parte del sistema tenga que conocer todos los detalles de cada modelo, LiteLLM ofrece un punto único y homogéneo de acceso.

% ============================================================
\section{Docker y orquestación de contenedores}
\label{sec:docker}
% ============================================================

Docker\footnote{\url{https://www.docker.com/}} y Docker Compose son la base del despliegue reproducible. Se utilizan para empaquetar todos los componentes y ponerlos en marcha como un conjunto coordinado.

\subsection{Ventajas de la contenerización}

\begin{itemize}
    \item \textbf{Reproducibilidad}: el mismo sistema puede desplegarse de forma controlada en distintos entornos.
    \item \textbf{Aislamiento}: cada servicio mantiene sus dependencias separadas.
    \item \textbf{Seguridad perimetral}: solo ciertos servicios se exponen al exterior.
    \item \textbf{Acceso a GPU cuando es necesario}: se reserva aceleración solo para los servicios que la necesitan.
\end{itemize}

En el contexto del TFG, esto también es una parte de la contribución: no solo se han elegido tecnologías, sino que se han dejado integradas dentro de una arquitectura desplegable de forma repetible.

% ============================================================
\section{Herramientas de procesamiento documental y web}
\label{sec:herramientas-procesamiento}
% ============================================================

Esta sección reúne herramientas ya existentes que no han sido desarrolladas en el proyecto, pero que sí se han integrado y combinado para cubrir necesidades concretas del sistema.

\subsection{PaddleOCR: reconocimiento óptico de caracteres}

PaddleOCR\footnote{\url{https://github.com/PaddlePaddle/PaddleOCR}} es un motor OCR externo utilizado para extraer texto de PDFs escaneados \cite{paddleocr2024docs}. No se ha implementado un OCR propio; el trabajo realizado aquí consiste en decidir cuándo usarlo, cómo integrarlo en la ingesta y cómo hacerlo convivir con el resto de componentes del sistema.

\subsection{Playwright: automatización de navegadores}

Playwright\footnote{\url{https://playwright.dev/}} es una herramienta externa de automatización de navegador \cite{playwright2024docs}. Se utiliza porque muchas páginas web modernas dependen de JavaScript y no pueden procesarse correctamente con una simple petición HTTP. Su papel dentro del proyecto es permitir obtener el contenido real que ve el usuario en páginas dinámicas.

\subsection{Trafilatura: extracción de contenido web}

Trafilatura es una biblioteca externa especializada en extraer el contenido principal de una página web, eliminando menús, barras laterales y ruido. En este TFG no se ha desarrollado un extractor propio de texto web desde cero; se ha integrado Trafilatura como primera opción dentro de una cadena de extracción con \textit{fallbacks}, porque resuelve bien el problema de aislar el contenido relevante.

\subsection{SentenceTransformers: generación de embeddings}

SentenceTransformers \cite{reimers2019sentence,sbertDocs} es una biblioteca externa utilizada para generar embeddings. Se emplea el modelo \texttt{paraphrase-multilingual-MiniLM-L12-v2}, elegido por su equilibrio entre tamaño, soporte multilingüe y facilidad de despliegue. Aquí también es importante la distinción metodológica: el modelo y la librería ya existían; lo que aporta el proyecto es su integración en el pipeline RAG y su ajuste al caso de uso documental.

% ============================================================
\section{Stack de monitorización}
\label{sec:stack-monitorizacion}
% ============================================================

Para operar un sistema de este tipo no basta con que funcione: también hay que entender cómo se comporta. Por ello se incorpora un stack de observabilidad formado por herramientas ya existentes, configuradas y conectadas dentro del proyecto.

\subsection{Prometheus}

Prometheus\footnote{\url{https://prometheus.io/}} recoge métricas temporales del backend y de otros servicios. Su utilidad en este TFG es permitir medir latencia, tasa de errores y patrones de uso en lugar de basar las decisiones de optimización en impresiones subjetivas.

\subsection{Grafana}

Grafana\footnote{\url{https://grafana.com/}} visualiza las métricas recogidas por Prometheus mediante paneles comprensibles. En el proyecto se ha utilizado para construir paneles específicos de backend, GPU y bases de datos. Esto no debe entenderse como un desarrollo de una herramienta nueva, sino como el diseño de la capa de observación necesaria para operar el sistema con criterio.

\subsection{NVIDIA DCGM Exporter}

NVIDIA DCGM Exporter\footnote{\url{https://github.com/NVIDIA/dcgm-exporter}} expone métricas de la GPU en formato compatible con Prometheus. Su presencia se justifica porque en un sistema que ejecuta modelos y OCR en local, la GPU es un recurso crítico y debe monitorizarse con el mismo nivel de detalle que el backend.

% ============================================================
\section{Otras tecnologías del ecosistema}
\label{sec:otras-tecnologias}
% ============================================================

\begin{itemize}
    \item \textbf{PostgreSQL}: almacena usuarios, conversaciones y metadatos de OpenWebUI.
    \item \textbf{Redis}: se emplea como caché de respuestas a través de LiteLLM.
    \item \textbf{MinIO}: conserva documentos binarios como respaldo y fuente de reindexación.
    \item \textbf{NGINX}: actúa como proxy inverso y punto único de entrada.
    \item \textbf{MSAL}: permite autenticación con Azure AD y acceso autenticado a SharePoint mediante Microsoft Graph.
\end{itemize}

% ============================================================
\section{Diagrama de flujo de tecnologías usadas}
\label{sec:diagrama-flujo}
% ============================================================

La Figura~\ref{fig:flujo-tecnologias} resume cómo se conectan las tecnologías principales del sistema. No representa únicamente el recorrido de una consulta, sino también los componentes de soporte que permiten que el sistema funcione en un entorno real: autenticación, almacenamiento, caché y monitorización. Se incluye aquí porque ayuda a que un lector no especializado entienda de un vistazo qué piezas existen y cómo se relacionan.

\begin{figure}[H]
\centering
\begin{tikzpicture}[
    node distance=0.9cm and 0.8cm,
    tech/.style={draw, rounded corners, align=center, minimum width=2.7cm, minimum height=0.9cm, fill=gray!10},
    own/.style={tech, fill=blue!8},
    ext/.style={tech, fill=green!8},
    infra/.style={tech, fill=orange!10},
    arrow/.style={-Latex, thick},
    area/.style={draw, dashed, rounded corners, inner sep=0.35cm}
]
\node[ext] (user) {Usuario};
\node[infra, below=of user] (nginx) {NGINX\\Proxy de entrada};
\node[ext, below left=1.0cm and 1.3cm of nginx] (azure) {Azure AD / MSAL\\Autenticación};
\node[ext, below=of nginx] (ui) {OpenWebUI\\Interfaz web};
\node[own, below=of ui] (pipeline) {Pipeline JARVIS\\Python};
\node[own, below left=1.4cm and 2.0cm of pipeline] (backend) {Backend RAG\\FastAPI};
\node[own, below right=1.4cm and 2.0cm of pipeline] (indexer) {Indexer\\FastAPI};
\node[own, right=of pipeline] (boe) {Servidor MCP /\\Integración BOE};

\node[ext, below=of backend] (qdrant) {Qdrant\\Base vectorial};
\node[ext, left=of qdrant] (litellm) {LiteLLM};
\node[ext, left=of litellm] (ollama) {Ollama\\LLMs locales};
\node[ext, right=of qdrant] (postgres) {PostgreSQL\\Usuarios y chats};
\node[ext, below=of qdrant] (redis) {Redis\\Caché};
\node[ext, right=of indexer] (sharepoint) {SharePoint\\Microsoft Graph};
\node[ext, below=of indexer] (ocr) {PaddleOCR\\OCR};
\node[ext, left=of indexer] (web) {Playwright +\\Trafilatura};
\node[ext, below=of redis] (minio) {MinIO\\Documentos binarios};

\node[infra, right=of postgres] (prom) {Prometheus\\Métricas};
\node[infra, right=of prom] (grafana) {Grafana\\Dashboards};
\node[infra, below=of prom] (dcgm) {NVIDIA DCGM\\GPU metrics};

\draw[arrow] (user) -- (nginx);
\draw[arrow] (nginx) -- (ui);
\draw[arrow] (ui) -- (azure);
\draw[arrow] (ui) -- (pipeline);
\draw[arrow] (pipeline) -- (backend);
\draw[arrow] (pipeline) -- (ui);
\draw[arrow] (pipeline) -- (indexer);
\draw[arrow] (pipeline) -- (boe);
\draw[arrow] (backend) -- (qdrant);
\draw[arrow] (backend) -- (litellm);
\draw[arrow] (backend) -- (postgres);
\draw[arrow] (backend) -- (redis);
\draw[arrow] (backend) -- (minio);
\draw[arrow] (litellm) -- (ollama);
\draw[arrow] (indexer) -- (sharepoint);
\draw[arrow] (indexer) -- (ocr);
\draw[arrow] (indexer) -- (web);
\draw[arrow] (indexer) -- (qdrant);
\draw[arrow] (indexer) -- (minio);
\draw[arrow] (backend) -- (prom);
\draw[arrow] (dcgm) -- (prom);
\draw[arrow] (prom) -- (grafana);

\node[area, fit=(nginx)(azure)(ui)(pipeline)] {};
\node[area, fit=(backend)(indexer)(boe)] {};
\node[area, fit=(ollama)(litellm)(qdrant)(postgres)(redis)(minio)] {};
\node[area, fit=(prom)(grafana)(dcgm)] {};
\end{tikzpicture}
\caption{Diagrama general de tecnologías usadas en el sistema. En azul se muestran los componentes desarrollados en el TFG, en verde las herramientas reutilizadas y en naranja la infraestructura de soporte y monitorización.}
\label{fig:flujo-tecnologias}
\end{figure}
